\documentclass[10pt,letterpaper,final]{report}
\usepackage[margin=.75in]{geometry}
\usepackage[utf8]{inputenc}
\usepackage{amsmath}
\usepackage{amsfonts}
\usepackage{amssymb}
\usepackage{amsthm}
\renewcommand\qedsymbol{$\blacksquare$}
\usepackage{enumerate}
\usepackage{enumitem}
\usepackage{hyperref}
\usepackage{pdfpages}
\usepackage{graphics}
\usepackage{graphicx}
\usepackage{tikz}
\usepackage{tikz-qtree}
\usepackage{forest}
\usepackage{float}
\usetikzlibrary{automata,arrows}
\usetikzlibrary{shapes.multipart, positioning}
\usepackage{mdframed}
\usepackage[
  backend=biber,
  style=numeric-comp,
  sorting=none
]{biblatex}
\addbibresource{references.bib}

\usepackage{minted}
\usemintedstyle{algol_nu}
\usepackage{xltabular}
\usepackage{pdflscape}
\usepackage{tcolorbox}
\usepackage{array}
\usepackage{todonotes}
\setuptodonotes{inline}  % render \todo notes as full-width blocks in the text flow, not margin boxes

\newcommand{\reason}[1]{&& \text{#1}}

% Based on template from COSC-417 at Towson University, originally authored by Marius Zimand

% To input a script, adapt the following:
% \inputminted[breaklines, breakanywhere, fontsize=\footnotesize, frame=single, fontsize=\footnotesize]{python}{script.py}

\begin{document}
\begin{center}
  \fbox{
    \vbox{
      \begin{flushleft}
        Jonathan Llovet \\  % authors' names
        EN.625.252.81 \\  %class
        2026-11-12 \\  % date
        Module 11 - Discussion 11 \\ % ID
      \end{flushleft}
      \center{\Large{\textbf{Linear Algebra}}}
      %\end{mdframed}
    } % end vbox
  } % end fbox
  \vline
\end{center}

\noindent Answer any \textbf{two} of the following.

\medskip
\noindent\textbf{Prompt 1:}
\medskip

The dot product of vectors is a specific example of a special function called an \emph{inner product} on a particular set $V$. Define an inner product in general, and give an example of an inner product on another set other than $R^n$.

\medskip
\noindent\textbf{Answer:}
\medskip

\pagebreak
\noindent\textbf{Prompt 2:}
\medskip

The norm of a vector, defined using the dot product, is an attempt to define the notion of distance in $R^n$. This norm is called the \emph{Euclidean norm}.  As with everything, the norm can be thought of as a function on vectors to the nonnegative real line that satisfies certain properties. Define the abstract norm or distance function $d(x,y)$.  Then find and define other norms on $R^n$. Some examples include the Taxicab norm, max norm, or $p$-norm.

\medskip
\noindent\textbf{Answer:}
\medskip

\pagebreak
\noindent\textbf{Prompt 3:}
\medskip

State the Cauchy-Schwarz inequality.  Then explain the following statement ``The triangle inequality is just an application of the Cauchy-Schwarz inequality.'' in your own words. Then define a \emph{Metric Space}, noting the use of the triangle inequality. Given an example of a Metric Space other than $R^n$.

\medskip
\noindent\textbf{Answer:}
\medskip

\pagebreak
\noindent\textbf{Prompt 4:}
\medskip

Define an \emph{orthogonal matrix}. Explain why the determinant of an orthogonal matrix must be $1$ or $-1$. Give some computational advantages or discuss some decompositions that using orthogonal matrices.

\medskip
\noindent\textbf{Answer:}
\medskip

\pagebreak

\section*{Source Code}
The full source code, including LaTeX, tooling, other artifacts, and git history is available on my Github repository for the class here:
\begin{center}
  \url{https://github.com/jllovet/jhu-en-625-252-fa26-linear-algebra}
\end{center}

\printbibliography

\end{document}
